\documentclass[11pt]{article}

\usepackage[T1]{fontenc}
\usepackage[utf8]{inputenc}
\usepackage{lmodern}
\usepackage{amsmath,amssymb,amsthm,mathtools}
\usepackage{microtype}
\usepackage[margin=1in]{geometry}
\usepackage{booktabs,longtable,array}
\usepackage{enumitem}
\usepackage{xcolor}
\usepackage{listings}
\usepackage{hyperref}
\usepackage{url}
\usepackage{fancyhdr}

\definecolor{deepblue}{RGB}{29,66,123}
\definecolor{softgray}{RGB}{245,246,248}
\definecolor{darkred}{RGB}{130,25,25}
\hypersetup{
  colorlinks=true,
  linkcolor=deepblue,
  urlcolor=deepblue,
  citecolor=deepblue,
  pdftitle={An Omega-Diagonal Forest Reflection Above BusyLean Contender 10},
  pdfauthor={Research proposal for BusyLean}
}

\pagestyle{fancy}
\fancyhf{}
\lhead{BusyLean Contender 11 proposal}
\rhead{October 2026}
\cfoot{\thepage}
\setlength{\headheight}{14pt}

\lstdefinelanguage{lean}{
  keywords={def,theorem,inductive,structure,where,match,with,fun,let,in,
    namespace,end,open,abbrev,import,by,have,show,exact,termination_by,
    unfold,change,simpa,using,cases,obtain,refine},
  sensitive=true,
  morecomment=[l]{--},
  morecomment=[s]{/-}{-/},
  morestring=[b]"
}
\lstset{
  language=lean,
  basicstyle=\ttfamily\footnotesize,
  keywordstyle=\bfseries\color{deepblue},
  commentstyle=\itshape\color{black!58},
  columns=fullflexible,
  keepspaces=true,
  frame=single,
  framerule=0.3pt,
  rulecolor=\color{black!25},
  backgroundcolor=\color{softgray},
  xleftmargin=0.8em,
  xrightmargin=0.8em,
  aboveskip=0.8em,
  belowskip=0.8em,
  breaklines=true,
  breakatwhitespace=false,
  literate=
    {→}{{$\to$}}1
    {←}{{$\leftarrow$}}1
    {ℕ}{{$\mathbb N$}}1
    {∀}{{$\forall$}}1
    {∃}{{$\exists$}}1
    {≤}{{$\le$}}1
    {<}{{$<$}}1
    {∈}{{$\in$}}1
    {×}{{$\times$}}1
    {λ}{{$\lambda$}}1
    {ω}{{$\omega$}}1
    {κ}{{$\kappa$}}1
    {·}{{$\cdot$}}1
    {₁}{{$_1$}}1
    {₂}{{$_2$}}1
    {₃}{{$_3$}}1
    {⟨}{{$\langle$}}1
    {⟩}{{$\rangle$}}1
}

\newtheorem{theorem}{Theorem}[section]
\newtheorem{lemma}[theorem]{Lemma}
\newtheorem{proposition}[theorem]{Proposition}
\newtheorem{corollary}[theorem]{Corollary}
\theoremstyle{definition}
\newtheorem{definition}[theorem]{Definition}
\newtheorem{construction}[theorem]{Construction}
\theoremstyle{remark}
\newtheorem{remark}[theorem]{Remark}
\newtheorem{warning}[theorem]{Verification warning}

\newcommand{\N}{\mathbb N}
\newcommand{\Fk}{\mathrm{F}_{\kappa}}
\newcommand{\WTree}{\mathsf W}
\newcommand{\Round}{\mathcal R}
\newcommand{\Tower}{\mathcal T}
\newcommand{\lean}[1]{\texttt{\detokenize{#1}}}
\newcommand{\Cten}{C_{10}}
\newcommand{\Celeven}{C_{11}}
\newcommand{\Bfam}{B}
\newcommand{\defeq}{\mathrel{:=}}
\newcommand{\code}[1]{\texttt{\detokenize{#1}}}

\title{\textbf{An Omega-Diagonal Forest Reflection\\Above BusyLean Contender 10}\\[0.5em]
  \large A computable Lean proposal with an explicit comparison architecture}
\author{Research proposal prepared for the BusyLean project}
\date{October 12, 2026}

\begin{document}
\maketitle

\begin{abstract}
BusyLean asks for closed computable natural numbers in Lean, together with
kernel-checkable proofs that each new entry exceeds its predecessor. The current
repository entry, Contender~10, reflects kind-polymorphic System~F$_\omega$
programs over Contender~9's family. We propose Contender~11, which combines that
language-level reflection with the repository's strongest existing structural
amplifier: forest reflection along the self-generated W-tree index that joins
every tree denoted by the W-language. One traversal is promoted to a function
operator $\Round$. We then define a self-accelerated tower
$\Tower_{0}=\Bfam$ and
$\Tower_{k+1}(n)=\Round(\Tower_k)(\Tower_k(n))$, where
$\Bfam(n)=F_{10}(F_{10}(n))$ is Contender~10's construction as a family. The
new diagonal is
\[
  F_{11}(n)=\Tower_{\Bfam(n)+1}(\Bfam(n)),
  \qquad
  \Celeven=F_{11}(F_{11}(42)).
\]
Every component is computable in principle. The comparison avoids any invalid
monotonicity assumption about oracle reflection: it exhibits the seed explicitly
inside the index tree, proves that every tower successor is larger by the
round's input bound, and derives
$\Cten < F_{11}(42)<\Celeven$. The package contains Lean definitions, proposed
proofs, a transitive dependency audit, a clean-build checker, and this article.
Proof-theoretic size claims are deliberately separated from the exact formal
inequality.
\end{abstract}

\begin{center}
\fcolorbox{darkred}{softgray}{\parbox{0.91\textwidth}{
\textbf{Verification status.} The article has been compiled and visually
checked. The Lean source has passed a conservative source-level audit, but the
execution environment used to prepare this package did not contain Lean~4.34.0.
Consequently, the modules in this archive are a serious, integration-ready
\emph{proposal}, not a claim of completed kernel verification. The bundled
\code{check11.py} and \code{Contender11Audit.lean} specify the exact final
checks.}}
\end{center}

\tableofcontents
\newpage

%======================================================================
\section{Problem, constraints, and baseline}

\subsection{The BusyLean game}

The public BusyLean repository is a Lean~4 continuation of Cody Roux's
``name the biggest number'' contest. A successful entry is not merely a concise
large-number notation. It must be a closed term of type \lean{Nat}, computable
in principle, and it must be accompanied by a formal strict inequality against
the preceding entry. The repository rules further prohibit escape hatches such
as custom axioms, \lean{noncomputable}, \lean{partial}, \lean{unsafe}, opaque
stand-ins, external implementations, and \lean{native_decide}. The only allowed
logical axioms are \lean{propext} and \lean{Quot.sound}; the current target is
Lean~4.34.0. Definition and proof modules are also subject to separate time
budgets. These constraints make the problem a useful synthesis of proof theory,
programming-language semantics, formal verification, and careful software
engineering \cite{BusyLean,Roux}.

The no-laziness rule matters. A successor may use the preceding
\emph{construction} as a seed or oracle for a qualitatively stronger operation,
but it may not simply mention the previous constant and add one, exponentiate
it, or wrap it in a fixed arithmetic expression. A convincing successor must
expose a new source of growth and prove that the predecessor occurs inside it.

\subsection{The current incumbent}

At the inspected repository revision, Contender~10 is

\begin{lstlisting}
def family9 (n : Nat) : Nat := F9 (F9 n)
def F10 (n : Nat) : Nat := Fk.reflectFk family9 n
def contender_10 : Nat := F10 (F10 42)
\end{lstlisting}

Here \lean{Fk.reflectFk} takes a finite maximum over accepted, closed programs
of kind-polymorphic System~F$_\omega$ at a syntax/fuel budget~$n$, evaluated
relative to an oracle~$f$, and also exceeds $f(n)$ and~$n$. The repository has a
strong-normalization theorem for the object language, so program evaluation is
a total Lean function rather than a bounded simulation or a semantic oracle.
The accepted comparison proves
\lean{contender_9_lt_contender_10} \cite{BusyLeanC10,BusyLeanComp10}.

Contender~10 is already far beyond elementary or primitive-recursive
large-number games. The problem is therefore not to add a familiar fast-growing
function. A meaningful successor should exploit a structural operation that is
not already compressed into the single bounded reflection step.

\subsection{Available machinery that has not yet been combined}

Two ingredients coexist in the repository:

\begin{enumerate}[leftmargin=*]
\item \textbf{Kind-polymorphic reflection.} Contender~10 can normalize and
  decode accepted programs of a language stronger than plain System~F$_\omega$,
  with a computable oracle.
\item \textbf{Self-generated forest reflection.} Contenders~6 and~9 use a
  well-founded, countably branching tree \lean{index}. Its children join all
  W-trees denoted by programs of the repository's W-language. At every node,
  the reflected program sees the family of child stages as an oracle, so it can
  select a child dynamically instead of losing information in a pointwise
  maximum \cite{BusyLeanC6,BusyLeanC9}.
\end{enumerate}

The repository's own roadmap observes that forest reflection can be applied to
new languages. The first move of the present proposal is precisely to apply the
Contender~9 amplifier to Contender~10's stronger reflected language. The second
move is to iterate the resulting complete operation and diagonalize over the
number of iterations.

%======================================================================
\section{Design criteria}

A serious successor should satisfy more than the final inequality. We impose the
following criteria.

\paragraph{C1: closed computability.}
The final declaration must reduce, in principle, to a natural number by Lean's
ordinary evaluator. All recursion must be structural or justified by an already
formalized well-founded argument.

\paragraph{C2: no predecessor constant.}
The definition may refer to \lean{F10} and the preceding construction as a
function, but its transitive declaration closure must not contain
\lean{contender_10} or any earlier contender constant.

\paragraph{C3: qualitative amplification.}
The new operation must do more than evaluate the old family at a larger fixed
argument. It should expand the reflected program class, the iteration
structure, or both.

\paragraph{C4: explicit comparison.}
Reflection relative to an oracle is generally not monotone: a program may branch
on an oracle value, so pointwise larger oracle answers can produce a smaller
result. Every inequality must therefore come from an explicit occurrence in a
finite maximum or from a child embedded in a known tree.

\paragraph{C5: modest proof surface.}
The construction should reuse the repository's difficult semantic theorem,
namely strong normalization for kind-polymorphic System~F$_\omega$, instead of
introducing a new object language whose metatheory cannot be completed within
the contest budgets.

\paragraph{C6: honest calibration.}
Lean will certify the strict numerical comparison. It will not automatically
certify an informal ordinal assignment. The article must distinguish exact
formal statements from proof-theoretic heuristics.

These criteria rule out a number of superficially attractive definitions. For
example, $2^{\Cten}$ is computable and larger, but it violates C2 and C3. A
fixed tenfold iterate of \lean{F10} uses the construction rather than the
constant, but still offers no new language or reflection structure. At the
other extreme, formalizing a Calculus of Inductive Constructions with large
elimination could provide a dramatic language-level jump, but violates C5 for a
compact successor proposal.

%======================================================================
\section{Reused formal ingredients}

\subsection{Countable oracle families}

A forest node has countably many children. The reflected language accepts one
unary oracle, so Contender~9 uses the computable coding
\[
  \langle i,m\rangle_2 = 2^i(2m+1).
\]
The decoder takes the $2$-adic valuation to recover~$i$ and then recovers~$m$.
The theorem \lean{unpair2_code} proves that the code is inverted exactly. Thus
for a family $fs:\N\to\N\to\N$,
\[
  \lean{packFamily}\;fs\;(2^i(2m+1))=fs(i,m).
\]
The coding is not itself the source of growth. Its role is information
preservation: a program can compute a child index and an argument at run time,
then call that exact child stage.

\subsection{The self-generated W-tree index}

The repository's W-language contains W-types over finite types and a total
interpreter. A program that denotes a tree of type $\WTree\N$ is converted to a
Lean tree by \lean{defTree}; a non-tree value defaults to a leaf. For each
syntax depth~$k$, \lean{treesOfDepth k} is a node whose children enumerate the
trees denoted at that depth. The global index is
\[
  \lean{index}=\lean{node}\;(k\mapsto\lean{treesOfDepth}\;k).
\]
It is therefore a well-founded tree in Lean, yet it joins every W-tree that the
object language denotes at some finite syntax depth. The key formal embedding
used repeatedly in BusyLean is
\[
  \text{leaf}
  \prec \lean{treesOfDepth 0}
  \prec \lean{index},
\]
where each $\prec$ is an actual parent-child relation visible at every
sufficient input. The first enumerated depth-zero tree is definitionally a leaf,
recorded by \lean{first_tree_leaf}.

\subsection{Kind-polymorphic program evaluation}

The namespace \lean{BusyLean.Fk} provides:

\begin{itemize}[leftmargin=*]
\item finite enumerations \lean{kindsUpTo}, \lean{tysUpTo}, and
  \lean{ctmsUpTo};
\item a sound fuel-bounded checker;
\item a strong-normalization theorem for accepted erased terms;
\item a total normalizer and Church-numeral decoder; and
\item \lean{fvalue f fuel t}, the numerical value of an accepted program
  relative to the unary oracle~$f$, with zero as the rejection/default value.
\end{itemize}

This is the expensive semantic core. The proposal treats it as a verified
component and builds a stronger reflection geometry around it. Reducibility
candidates are the classical technique behind strong-normalization arguments for
polymorphic calculi; related modular normalization results explain why language
boundaries in the lambda cube are central to proof-theoretic comparisons
\cite{GeuversNederhof}.

%======================================================================
\section{The construction}

\subsection{One family-reflection step}

Let $fs:\N\to\N\to\N$ be a countable family. Define
\begin{align*}
M_{fs}(n)
  &= \max\left(
       \max_{i\le n}fs(i,n),
       \max_{t\in\mathrm{CTm}_{\le n}}
         \mathrm{value}_{\Fk}(\mathrm{pack}(fs),n,t)
     \right),\\
\mathrm{reflect}_{11}(fs,n)
  &= n+1+M_{fs}(n).
\end{align*}
The Lean definition is:

\begin{lstlisting}
def reflect11 (fs : W.Family) (n : Nat) : Nat :=
  n + 1 + BigNumber.maxN
    (W.prefixMax (fun i => fs i n) n)
    (BigNumber.maximum
      (Fk.fvalue (packFamily fs) n)
      (Fk.ctmsUpTo n))
\end{lstlisting}

The added $n+1$ is not intended as a growth amplifier. It provides a robust
strict inequality above the input and above either maximum, simplifying every
subsequent comparison.

\subsection{Forest reflection}

The family step is iterated structurally over a W-tree:

\begin{lstlisting}
def stage11 (seed : W.Family) : BigNumber.WTree Nat → Nat → Nat
  | .leaf   => reflect11 seed
  | .node g => reflect11 (fun i => stage11 seed (g i))
\end{lstlisting}

At a leaf, the reflected language receives the seed family. At a node, it
receives the family of child stages. Notice what is \emph{not} done: the
children are not collapsed to a pointwise maximum before program evaluation.
The program sees them through \lean{packFamily}, so it can choose a child based
on intermediate computation.

\subsection{A complete round}

For a unary seed $f:\N\to\N$, repeat it across the leaf family and traverse the
full self-generated index:

\begin{lstlisting}
def round11 (f : Nat → Nat) (n : Nat) : Nat :=
  stage11 (fun _ => f) index n
\end{lstlisting}

We write this function operator as $\Round(f)$. A complete round is already a
substantial candidate successor: it reflects the kind-polymorphic language at
every node along every W-language-definable tree embedded in \lean{index}.
However, the operation itself is computable and strictly seed-increasing, so it
can be iterated.

\subsection{The predecessor family}

The rule against direct use of the old constant is respected by exposing the
entire preceding construction as a function:

\begin{lstlisting}
def family10 (n : Nat) : Nat := F10 (F10 n)
\end{lstlisting}

Write $\Bfam(n)=F_{10}(F_{10}(n))$. Then $\Bfam(42)$ is definitionally equal to
Contender~10, but the declaration \lean{contender_10} itself does not occur in
the new definition.

\subsection{Self-accelerated iteration}

A naive iteration would set
$S_{k+1}(n)=\Round(S_k)(n)$. We use a stronger and still elementary recursion:
\[
  \Tower_0(n)=\Bfam(n),
  \qquad
  \Tower_{k+1}(n)=\Round(\Tower_k)(\Tower_k(n)).
\]
Thus the preceding stage plays two roles. It is the oracle seed throughout the
next forest, and its value becomes the next syntax/fuel/input budget. The Lean
code is:

\begin{lstlisting}
def tower11 : Nat → Nat → Nat
  | 0     => family10
  | k + 1 => fun n => round11 (tower11 k) (tower11 k n)
\end{lstlisting}

This recursion is structural on~$k$. It does not attempt to evaluate the huge
natural during type checking. As a mathematical function it is total for every
finite~$k$ and~$n$.

\subsection{The diagonal and final number}

The number of complete rounds should not be a small numeral chosen by the
paper's author. We take it from the preceding construction itself:

\begin{lstlisting}
def F11 (n : Nat) : Nat :=
  tower11 (family10 n + 1) (family10 n)

def contender_11 : Nat := F11 (F11 42)
\end{lstlisting}

Equivalently,
\[
  F_{11}(n)=\Tower_{\Bfam(n)+1}(\Bfam(n)),
  \qquad
  \Celeven=F_{11}(F_{11}(42)).
\]
The family $\Bfam$ is provably larger than its input, so the iteration depths
$\Bfam(n)+1$ are unbounded. In this precise but modest sense, $F_{11}$ is an
$\omega$-diagonal through all finite numbers of complete rounds. It is not a
literal transfinite recursion of length~$\omega$; every evaluation uses a finite
but extraordinarily large number of stages.

%======================================================================
\section{Exact comparison architecture}

This section gives the mathematical content mirrored by
\lean{Comparison11.lean}. The difficult metatheory is already inside
\lean{Fk.fvalue}; the new comparison is elementary arithmetic plus explicit
child inclusions.

\subsection{Local maximum inequalities}

\begin{lemma}[Visible oracle values]
If $i\le n$, then
\[
  fs(i,n)<\mathrm{reflect}_{11}(fs,n).
\]
\end{lemma}

\begin{proof}
The prefix maximum contains $fs(i,n)$. It is at most the outer maximum, which
is strictly below $n+1$ plus that maximum.
\end{proof}

\begin{lemma}[Accepted program values]
If $t$ belongs to \lean{Fk.ctmsUpTo n}, then
\[
 \lean{Fk.fvalue}\;(\lean{packFamily}\;fs)\;n\;t
 <\mathrm{reflect}_{11}(fs,n).
\]
\end{lemma}

\begin{proof}
The finite list maximum contains the displayed value, and the same arithmetic
argument applies.
\end{proof}

\begin{lemma}[Input bound]
For every $fs$ and~$n$,
\[
  n<\mathrm{reflect}_{11}(fs,n).
\]
\end{lemma}

The Lean proofs reuse the repository's lemmas for prefix and list maxima. No
normalization or evaluation of a contender is required.

\subsection{Children are below parents}

At a node $\lean{node}\;g$, the oracle family is
$i\mapsto\lean{stage11}\;seed\;(g(i))$. Therefore, when $i\le n$, the visible
oracle lemma yields
\[
  \mathrm{stage}_{11}(seed,g(i),n)
  <\mathrm{stage}_{11}(seed,\mathrm{node}(g),n).
\]
This is \lean{stage11_child_lt}. The input bound similarly lifts from one
reflection step to every tree.

\subsection{A full round explicitly contains its seed}

The critical proof is not a monotonicity theorem.

\begin{theorem}[Seed inclusion]
For every $f:\N\to\N$ and every~$n$,
\[
  f(n)<\Round(f)(n).
\]
\end{theorem}

\begin{proof}
Use three strict inequalities.

First, at the leaf the constant family $i\mapsto f$ visibly contains $f(n)$:
\[
  f(n)<\mathrm{stage}_{11}(i\mapsto f,\mathrm{leaf},n).
\]
Second, the first child of \lean{treesOfDepth 0} is definitionally that leaf, so
child domination gives
\[
 \mathrm{stage}_{11}(i\mapsto f,\mathrm{leaf},n)
 <\mathrm{stage}_{11}(i\mapsto f,\lean{treesOfDepth 0},n).
\]
Third, \lean{treesOfDepth 0} is the zeroth child of \lean{index}, hence
\[
 \mathrm{stage}_{11}(i\mapsto f,\lean{treesOfDepth 0},n)
 <\mathrm{stage}_{11}(i\mapsto f,\lean{index},n)=\Round(f)(n).
\]
Transitivity completes the proof.
\end{proof}

The corresponding Lean theorem is:

\begin{lstlisting}
theorem seed_lt_round11 (f : Nat → Nat) (n : Nat) :
    f n < round11 f n := by
  have h₁ : f n < stage11 (fun _ => f) .leaf n := ...
  have h₂ : stage11 (fun _ => f)
      (defTree (nthD (BigNumber.allTerms 0) 0 .zero)) n <
      stage11 (fun _ => f) (treesOfDepth 0) n := ...
  have h₃ : stage11 (fun _ => f) (treesOfDepth 0) n <
      stage11 (fun _ => f) index n := ...
  rw [first_tree_leaf] at h₂
  exact Nat.lt_trans h₁ (Nat.lt_trans h₂ h₃)
\end{lstlisting}

This explicit path is robust under nonmonotone oracle behavior. A reflected
program may do anything with other oracle entries; the finite maximum still
contains the particular child needed for the proof.

\subsection{Every tower successor is larger}

\begin{theorem}[Strict tower]
For all $k,n\in\N$,
\[
  \Tower_k(n)<\Tower_{k+1}(n).
\]
\end{theorem}

\begin{proof}
By definition,
\[
 \Tower_{k+1}(n)=\Round(\Tower_k)(\Tower_k(n)).
\]
Every round exceeds the numerical input at which it is evaluated, so
\[
 \Tower_k(n)<\Round(\Tower_k)(\Tower_k(n)).
\]
\end{proof}

It follows by ordinary induction and transitivity that every positive tower
level exceeds the base family at the same input:
\[
  \Bfam(n)<\Tower_{k+1}(n).
\]
This result is formalized as \lean{family10_lt_tower11_succ}.

\subsection{The predecessor family exceeds its input}

Contender~10's local input theorem says
$n<F_{10}(n)$. Applying it twice gives
\[
  n<F_{10}(n)<F_{10}(F_{10}(n))=\Bfam(n).
\]
Thus \lean{input_lt_family10} requires no monotonicity of $F_{10}$: it uses
the input theorem once at~$n$ and once at the concrete input $F_{10}(n)$.

\subsection{The diagonal dominates the predecessor family}

\begin{theorem}[Family domination]
For every~$n$,
\[
  \Bfam(n)<F_{11}(n).
\]
\end{theorem}

\begin{proof}
First apply the preceding input bound at input $\Bfam(n)$:
\[
  \Bfam(n)<\Bfam(\Bfam(n)).
\]
The right-hand side is the base level
$\Tower_0(\Bfam(n))$. Since $\Bfam(n)+1$ is positive, strict tower iteration
then yields
\[
  \Bfam(\Bfam(n))
  <\Tower_{\Bfam(n)+1}(\Bfam(n))=F_{11}(n).
\]
\end{proof}

Combining this theorem with $n<\Bfam(n)$ gives the useful diagonal bound
$n<F_{11}(n)$.

\subsection{The required comparison}

At~$42$, $\Bfam(42)$ is definitionally Contender~10. Therefore
\[
  \Cten=\Bfam(42)<F_{11}(42).
\]
The diagonal input theorem, now at the concrete input $F_{11}(42)$, gives
\[
  F_{11}(42)<F_{11}(F_{11}(42))=\Celeven.
\]
Hence:

\begin{theorem}[Proposed Lean theorem]
\[
  \Cten<\Celeven.
\]
In the package this is the declaration
\lean{contender_10_lt_contender_11}.
\end{theorem}

The proof term never needs to normalize either number. It compares their
construction trees symbolically.

%======================================================================
\section{Structural margins beyond the required inequality}

A contest entry is easier to assess when it proves more than one isolated
comparison. The package includes three structural margins.

\subsection{Every round beats the seed pointwise}

The theorem \lean{seed_lt_round11} holds for an arbitrary unary function, not
just \lean{family10}. Thus $\Round$ is inflationary on the diagonal values
needed by the construction, even though no general pointwise monotonicity claim
is made.

\subsection{Every W-definable tree is eventually absorbed}

Suppose a W-language term~$t$ occurs in the finite enumeration at depth~$k$.
The repository lemma \lean{exists_nthD_of_mem} supplies an index~$i$ at which
$t$ occurs. For every $n\ge k+i$, both~$k$ and~$i$ are visible child indices.
Two applications of child domination establish
\[
 \mathrm{stage}_{11}(i\mapsto f,\mathrm{defTree}(t),n)
 <\Round(f)(n).
\]
This is formalized as \lean{defined_tree_eventually_lt11}. It records the
reason \lean{index} is stronger than a single hand-written tree: every tree
denoted by the object language is eventually present.

\subsection{Every visible accepted kind-polymorphic program is absorbed}

At the root of \lean{index}, the child family is
\[
 i\longmapsto
 \mathrm{stage}_{11}(j\mapsto f,\lean{treesOfDepth}(i)).
\]
For every accepted term in \lean{Fk.ctmsUpTo n}, the program maximum inside the
root reflection contains its value relative to this exact family. The theorem
\lean{fvalue_lt_round11} exposes that fact. It is the language-level analogue
of defined-tree absorption.

Together these statements show that the proposal is not merely a numerically
larger wrapper. A round absorbs both the relevant program enumeration and the
relevant tree enumeration, while preserving dynamic access to the family of
sub-stages.

%======================================================================
\section{Computability and totality}

The phrase ``computable in principle'' needs a concrete decomposition.

\subsection{Finite syntax enumeration}

For each natural budget~$n$, \lean{Fk.ctmsUpTo n} is a finite list generated by
primitive recursion. The list may be astronomically large, but finiteness, not
feasibility, is the contest requirement.

\subsection{Total program evaluation}

For each listed term, the checker either rejects it, returning~$0$, or supplies
a typing derivation at the Church-numeral type. The repository's reducibility
argument proves strong normalization of the erasure relative to every total
oracle function. The normalizer is then a well-founded recursive function and
the decoder is total. No halting oracle or Busy Beaver occurs.

\subsection{Finite maxima}

The maximum of a finite list of natural values is computable. Prefix maxima over
indices $0,\ldots,n$ are primitive recursive.

\subsection{Structural recursion over W-trees}

The function \lean{stage11} is defined by the recursor of the inductive W-tree
type. Although a node has a function of countably many children, every recursive
call is on an immediate subtree, and Lean's structural recursion principle
certifies totality.

\subsection{Finite tower recursion}

For each natural~$k$, \lean{tower11 k} is obtained by $k$ structural successor
steps. The fact that $k=\Bfam(n)+1$ is too large to realize physically does not
make it noncomputable. The recursion is finite for every input.

\subsection{Closedness}

All parameters are bound before the final declaration. The numeral~$42$, the
existing computable functions, the syntax enumerators, and the total
normalizers determine one closed \lean{Nat}. There are no external files,
random choices, unspecified witnesses, or semantic oracles in its definition.

\begin{remark}
Practical evaluability is intentionally absent. Even very small budgets for the
underlying syntax enumerations can be costly. BusyLean's rule is constructive
definability plus formal comparison, not decimal expansion or execution at the
contest input.
\end{remark}

%======================================================================
\section{Why the growth jump is significant}

\subsection{Three nested amplifiers}

The construction combines three different operations rather than repeatedly
applying one arithmetic function.

\begin{enumerate}[leftmargin=*]
\item \textbf{Language reflection.} At each node, take a maximum over accepted
  kind-polymorphic System~F$_\omega$ programs with an oracle.
\item \textbf{Tree-indexed reflection.} Arrange these node reflections along a
  self-generated W-tree that contains every W-language-definable tree.
\item \textbf{Self-accelerated finite iteration.} Feed the entire preceding
  round back as the next oracle and use its current value as the next budget;
  then choose the iteration count from the preceding contender family.
\end{enumerate}

Contender~10 uses the first operation once. A single \lean{round11} combines the
first two. The tower repeatedly applies the combined operation, and $F_{11}$
diagonalizes through a predecessor-sized number of repetitions.

\subsection{A cautious operator heuristic}

Let $\Delta$ denote, only heuristically, the proof-theoretic increment associated
with one complete $\Fk$ forest-reflection round over \lean{index}. If ordinary
oracle reflection behaved like ordinal addition, then repeated rounds would
suggest a progression
\[
  \beta,
  \quad \beta+\Delta,
  \quad \beta+\Delta\cdot 2,
  \quad \ldots,
\]
and a diagonal through unbounded finite counts would point toward a limit near
$\beta+\Delta\cdot\omega$. The self-feeding input
$\Tower_k(n)$ makes the numerical growth substantially more aggressive than
this schematic index progression alone.

This picture is useful intuition, but it is not a theorem for the implemented
operator. The reflected program can inspect exact oracle values, and the
self-generated tree index mixes syntax and semantic evaluation. A defensible
ordinal calibration would require a separate interpretation theorem connecting
the Lean hierarchy to a recognized reflection calculus or ordinal notation
system.

Modern proof theory confirms in many settings that iterated reflection is not a
cosmetic operation: reflection ranks can recover proof-theoretic ordinals, and
iterations of uniform reflection produce systematic ordinal progressions
\cite{PakhomovWalshRank,PakhomovWalshOmega}. Those results motivate the design,
but they do not identify the exact ordinal of this concrete BusyLean operator.

\subsection{Formal lower bounds that do not rely on the heuristic}

Even without an ordinal assignment, the package proves the following exact
facts, subject to kernel verification:

\begin{itemize}[leftmargin=*]
\item $\Bfam(n)<F_{11}(n)$ for every input~$n$;
\item $\Tower_k(n)<\Tower_{k+1}(n)$ for every finite level;
\item every W-language-definable tree stage is eventually below a round;
\item every visible accepted kind-polymorphic program value is below a round;
  and
\item $\Cten<\Celeven$.
\end{itemize}

These are sufficient to justify ``significantly exceeds'' in the contest sense:
the new number arises from a strict hierarchy of complete reflection rounds,
not from a one-off arithmetic padding term.

%======================================================================
\section{Why monotonicity is deliberately avoided}

It is tempting to assert
\[
  f\le g\quad\Longrightarrow\quad\Round(f)\le\Round(g).
\]
For an arbitrary reflected programming language this is false without strong
restrictions. A program can query the oracle and return~$100$ when the answer is
zero but return~$0$ otherwise. Increasing the oracle can therefore lower that
program's output. Taking a maximum over many programs does not automatically
repair the issue at every finite budget.

The comparison proofs use only four safe mechanisms:

\begin{enumerate}[leftmargin=*]
\item membership in a concrete finite maximum;
\item membership in a concrete prefix maximum;
\item a concrete parent-child edge in a W-tree; and
\item transitivity of strict order on natural numbers.
\end{enumerate}

This proof style is one of the proposal's strongest engineering properties. It
localizes every inequality and makes the argument insensitive to complicated
control flow inside the reflected language.

%======================================================================
\section{Lean engineering and audit plan}

\subsection{Module split}

The proposed definition module imports only \lean{BusyLean.Contender10}. It
contains the six computable declarations
\lean{reflect11}, \lean{stage11}, \lean{round11}, \lean{family10},
\lean{tower11}, \lean{F11}, and \lean{contender_11}. The comparison module
imports the definition and \lean{BusyLean.Comparison10}; this keeps proof-only
dependencies out of the number's definition module.

\subsection{Forbidden-feature source audit}

The bundled \code{static_audit.py} removes nested Lean comments and string
contents, then rejects occurrences of
\lean{sorry}, \lean{admit}, custom \lean{axiom}, \lean{noncomputable},
\lean{partial}, \lean{unsafe}, \lean{opaque}, \lean{implemented_by},
\lean{extern}, \lean{native_decide}, and \lean{Classical.choice}. It also
rejects any direct earlier contender name in the executable text of the
definition module and enforces a 100-column source limit.

This check is intentionally conservative and intentionally insufficient. Text
search cannot establish type correctness, termination, or transitive axiom
usage.

\subsection{Kernel-level dependency closure}

\code{Contender11Audit.lean} copies the repository's environment traversal. For
a root declaration it follows every constant in declaration types and bodies,
then fails if any forbidden earlier contender constant appears. The intended
assertion is:

\begin{lstlisting}
#assert_not_in_closure11 BusyLean.contender_11 :
  BusyLean.contender_10 BusyLean.contender_9 BusyLean.contender_8
  BusyLean.contender_7 BusyLean.contender_6 BusyLean.contender_5
\end{lstlisting}

This is stronger than checking the source file for a name. It detects an
indirect dependency introduced through helper declarations.

\subsection{Axiom reports}

The audit prints axioms for the new number, required comparison, tower theorem,
family theorem, tree absorption theorem, and program absorption theorem. The
checker parses Lean's output and accepts only \lean{propext} and
\lean{Quot.sound}, matching the current repository policy.

\subsection{Time budgets}

The proposed \code{check11.py} builds \lean{BusyLean.Comparison10} first as the
accepted predecessor, then times \lean{BusyLean.Contender11} against the
15-second definition budget and \lean{BusyLean.Comparison11} against the
60-second proof budget. It finally runs the audit. Actual timings must be
recorded on a reasonable machine before acceptance; none are fabricated here.

\subsection{Likely elaboration risks}

The mathematical proof uses only existing theorem interfaces, but this package
was prepared without a local Lean executable. The likely failure modes are
small elaboration issues rather than conceptual gaps:

\begin{itemize}[leftmargin=*]
\item definitional reduction may need an explicit \lean{simpa [F11]};
\item the root reduction in \lean{fvalue_lt_round11} may need
  \lean{unfold round11 index stage11};
\item the final comparison may prefer an explicit \lean{change}; or
\item a theorem name exported transitively by \lean{Comparison10} may need an
  additional direct import.
\end{itemize}

The package's \code{PROOF_OBLIGATIONS.md} provides fallback proof forms for
these points. None changes the construction or argument.

%======================================================================
\section{Alternatives considered}

\subsection{A single complete round}

The minimal successor
\[
  C' = \Round(\Bfam)(\Round(\Bfam)(42))
\]
would be easy to compare and already add forest reflection to Contender~10. It
was rejected as insufficiently ambitious because $\Round$ is itself a total,
strictly seed-increasing operator. Once that lemma is available, a finite tower
and diagonal cost little additional proof complexity.

\subsection{Same-input iteration}

One could define
$S_{k+1}(n)=\Round(S_k)(n)$. The selected recurrence evaluates at $S_k(n)$
instead. The stronger recurrence has an equally simple proof because every
round exceeds its input. It also makes the syntax budget, type-checker fuel, and
visible child prefix grow with the previous stage.

\subsection{Pointwise maxima of child stages}

Replacing the oracle family by
$\max_{i\le n}\mathrm{child}_i(n)$ would simplify types but destroy dynamic
child selection. A program would learn only one aggregate value, not the
identity and behavior of a selected child function. The existing packed-family
interface is strictly more informative.

\subsection{A universe ladder}

The BusyLean roadmap proposes stratified universes of kinds above the current
kind-polymorphic language. That is the most natural next language-level
research program. A complete entry would require syntax, substitution,
kind-checking, conversion, candidates indexed by universe level, adequacy, a
sound checker, and a normalizer. It may eventually dominate the present
amplifier by a qualitative margin, but it is not a low-risk continuation of the
current formalization.

\subsection{Dependent large elimination}

A CIC-like language with large elimination can let run-time naturals control
type structure and may surpass fixed finite-order hierarchies more directly.
Its realizability or normalization proof is substantially more involved. The
present proposal should be viewed as an immediately available successor and a
platform from which that larger project can be funded by a still larger oracle
family.

%======================================================================
\section{Integration procedure}

The archive is arranged so it can be copied into a BusyLean checkout without
replacing existing files.

\begin{enumerate}[leftmargin=*]
\item Copy \code{Contender11.lean} and \code{Comparison11.lean} into
  \code{BusyLean/}.
\item Copy \code{Contender11Audit.lean} into \code{test/}.
\item Copy \code{check11.py} into \code{scripts/}.
\item Append imports of the two new modules to \code{BusyLean.lean}.
\item Run \code{python3 scripts/check11.py} from the repository root.
\item If the build passes, inspect the axiom reports and closure message rather
  than relying only on the script's final line.
\item Add the supplied README section and update the main article only after
  timings and kernel verification are known.
\end{enumerate}

A unified patch is included for convenience. Because repository development may
continue, applying the patch with \code{git apply --check} before mutation is
recommended.

%======================================================================
\section{Further research}

\subsection{Formalize an abstract reflection-round interface}

The comparison pattern is independent of the details of~$\Fk$. An abstract
structure could package:

\begin{itemize}[leftmargin=*]
\item a finite program enumeration;
\item a total evaluator relative to an oracle family;
\item local maximum lemmas;
\item a forest-stage operator; and
\item a self-generated index with an explicit seed path.
\end{itemize}

Then seed inclusion, tower strictness, and diagonal comparison could be proved
once and instantiated for System~F, F$_\omega$, $\Fk$, and future languages.
This would reduce duplication and make future contenders easier to audit.

\subsection{Prove quantitative embedding theorems}

The current structural theorems show eventual domination of each individual
W-definable tree. A stronger result would assign each enumerated tree a concrete
level or path code and give an explicit threshold. This could support machine-
checked lower bounds against named fast-growing hierarchies.

\subsection{Study monotone sublanguages}

Although the full reflected language is nonmonotone in its oracle, a syntactic
positive-oracle fragment may admit a monotonicity theorem. Reflecting that
fragment could enable cleaner ordinal iteration arguments while retaining a
separate full-language maximum for raw size.

\subsection{Calibrate the round operator}

The main theoretical open problem is to identify an ordinal or reflection
calculus faithfully represented by $\Round$. Possible approaches include:

\begin{enumerate}[leftmargin=*]
\item interpreting each W-tree stage as an iterated reflection theory;
\item extracting a proof-theoretic dilator from the syntax-to-tree operation;
\item proving upper and lower simulations between the Lean hierarchy and a
  standard collapsing-function system; and
\item determining the effect of self-feeding the numeric budget on the
  function hierarchy at a fixed ordinal index.
\end{enumerate}

The literature on reflection ranks and reductions of model reflection to
iterated syntactic reflection offers conceptual tools, but the BusyLean object
languages and oracle-sensitive maxima require a bespoke theorem
\cite{PakhomovWalshRank,PakhomovWalshOmega,Rathjen}.

\subsection{Build the stratified kind-universe successor}

A major follow-up should formalize a predicative ladder
$\Box_0:\Box_1:\Box_2:\cdots$ above the impredicative term sort, with no
self-typing universe. The research tasks are:

\begin{itemize}[leftmargin=*]
\item syntax and de Bruijn operations for level-indexed kinds;
\item a sound checker with fuel-bounded conversion;
\item semantic domains and candidates indexed by levels;
\item renaming and substitution across universe contexts;
\item adequacy and strong normalization;
\item a total evaluator and reflection operator; and
\item comparison against Contender~11, preferably using \lean{F11} as oracle.
\end{itemize}

\subsection{Mechanize the metatheory of the audit itself}

The closure command is a metaprogram that inspects the environment. One can go
further by generating a signed manifest of declaration names, source hashes,
Lean version, build timings, and axiom reports. This would make contest entries
reproducible artifacts rather than informal snapshots.

%======================================================================
\section{Conclusion}

Contender~11 is designed as the strongest low-risk successor available from the
current BusyLean components. It does not attempt to win by arithmetic padding.
Instead, it lifts the incumbent's kind-polymorphic reflection to a family-aware
forest stage, traverses the self-generated index of all W-language-definable
trees, iterates complete rounds with self-fed budgets, and diagonalizes through
a number of rounds supplied by the predecessor construction itself.

The exact mathematical comparison is short because the architecture was chosen
for explicit inclusions:
\[
  \Cten=\Bfam(42)
  <F_{11}(42)
  <F_{11}(F_{11}(42))=\Celeven.
\]
No monotonicity of oracle reflection is assumed. Every component remains
computable in principle, and the definition avoids all earlier contender
constants. The remaining gate is empirical and formal rather than conceptual:
run Lean~4.34.0, resolve any minor elaboration differences, verify the axiom and
dependency audits, and record the build budgets. Until that is done, this is an
integration-ready research proposal; after it is done, it is a natural
Contender~11 for the repository.

%======================================================================
\appendix
\section{Core Lean definition}

For convenient review, the complete proposed definition module is reproduced
below.

\begin{lstlisting}
import BusyLean.Contender10

namespace BusyLean

def reflect11 (fs : W.Family) (n : Nat) : Nat :=
  n + 1 + BigNumber.maxN
    (W.prefixMax (fun i => fs i n) n)
    (BigNumber.maximum
      (Fk.fvalue (packFamily fs) n)
      (Fk.ctmsUpTo n))

def stage11 (seed : W.Family) : BigNumber.WTree Nat → Nat → Nat
  | .leaf   => reflect11 seed
  | .node g => reflect11 (fun i => stage11 seed (g i))

def round11 (f : Nat → Nat) (n : Nat) : Nat :=
  stage11 (fun _ => f) index n

def family10 (n : Nat) : Nat := F10 (F10 n)

def tower11 : Nat → Nat → Nat
  | 0     => family10
  | k + 1 => fun n => round11 (tower11 k) (tower11 k n)

def F11 (n : Nat) : Nat :=
  tower11 (family10 n + 1) (family10 n)

def contender_11 : Nat := F11 (F11 42)

end BusyLean
\end{lstlisting}

\section{Theorem dependency graph}

\begin{center}
\begin{tabular}{>{\raggedright\arraybackslash}p{0.34\textwidth}
                >{\raggedright\arraybackslash}p{0.57\textwidth}}
\toprule
Declaration & Main dependencies and role\\
\midrule
\lean{lt_reflect11}
  & arithmetic form of the maximum\\
\lean{oracle_lt_reflect11}
  & prefix membership $+$ outer maximum\\
\lean{fvalue_lt_reflect11}
  & list-maximum membership $+$ outer maximum\\
\lean{stage11_child_lt}
  & visible oracle entry at a parent node\\
\lean{seed_lt_round11}
  & leaf $\to$ depth-zero join $\to$ global index\\
\lean{tower11_step_lt}
  & round input bound at input $\Tower_k(n)$\\
\lean{family10_lt_tower11_succ}
  & induction and transitivity over tower levels\\
\lean{input_lt_family10}
  & two applications of Contender 10's input theorem\\
\lean{family10_lt_F11}
  & base-family growth plus positive tower level\\
\lean{input_lt_F11}
  & $n<\Bfam(n)<F_{11}(n)$\\
\lean{contender_10_lt_contender_11}
  & family domination at 42 plus the outer diagonal\\
\bottomrule
\end{tabular}
\end{center}

\section{Reproducibility commands}

\begin{lstlisting}[language={}]
# In the BusyLean repository root, after copying the package files:
python3 scripts/check11.py

# Individual checks:
lake build BusyLean.Contender11
lake build BusyLean.Comparison11
lake env lean test/Contender11Audit.lean

# Article build from this package:
cd article
latexmk -pdf -interaction=nonstopmode -halt-on-error contender11.tex
\end{lstlisting}

\section{Acceptance checklist}

\begin{enumerate}[leftmargin=*]
\item Both Lean modules elaborate under exactly Lean~4.34.0.
\item \lean{contender_11} is accepted as a computable \lean{Nat} definition.
\item The required strict comparison is kernel-checked without \lean{sorry}.
\item Every printed axiom set is a subset of
  $\{\lean{propext},\lean{Quot.sound}\}$.
\item The transitive closure of \lean{contender_11} contains no earlier
  contender constant.
\item Definition time is at most 15 seconds under the repository protocol.
\item Proof time is at most 60 seconds under the repository protocol.
\item The main README and article state measured timings, not estimates.
\item No one attempts to evaluate the contender.
\end{enumerate}

%======================================================================
\begin{thebibliography}{99}

\bibitem{BusyLean}
V. Reshetnikov and contributors.
\newblock \emph{BusyLean: name the biggest number, in Lean}.
\newblock \url{https://github.com/VladimirReshetnikov/BusyLean}, inspected
October 2026.

\bibitem{BusyLeanC6}
BusyLean contributors.
\newblock \emph{Contender 6 and forest reflection modules}.
\newblock Files \code{BusyLean/Reflection.lean},
\code{BusyLean/Contender6.lean}, and \code{BusyLean/Comparison.lean} in
\cite{BusyLean}.

\bibitem{BusyLeanC9}
BusyLean contributors.
\newblock \emph{Contender 9: System F$_\omega$ reflected along the
W-language's own trees}.
\newblock Files \code{BusyLean/Contender9.lean} and
\code{BusyLean/Comparison9.lean} in \cite{BusyLean}.

\bibitem{BusyLeanC10}
BusyLean contributors.
\newblock \emph{Contender 10: kind polymorphism}.
\newblock File \code{BusyLean/Contender10.lean} in \cite{BusyLean}.

\bibitem{BusyLeanComp10}
BusyLean contributors.
\newblock \emph{Contender 10 comparison theorems}.
\newblock File \code{BusyLean/Comparison10.lean} in \cite{BusyLean}.

\bibitem{Roux}
C. Roux and contributors.
\newblock \emph{name-the-biggest-number}.
\newblock \url{https://github.com/codyroux/name-the-biggest-number}.

\bibitem{Aaronson}
S. Aaronson.
\newblock Who can name the bigger number?
\newblock \url{https://www.scottaaronson.com/writings/bignumbers.html}.

\bibitem{GeuversNederhof}
H. Geuvers and M.-J. Nederhof.
\newblock Modular proof of strong normalization for the Calculus of
Constructions.
\newblock \emph{Journal of Functional Programming}, 1(2):155--189, 1991.
\newblock doi:10.1017/S0956796800020037.

\bibitem{PakhomovWalshRank}
F. Pakhomov and J. Walsh.
\newblock Reflection ranks and ordinal analysis.
\newblock \emph{Journal of Symbolic Logic}, 86(4):1350--1384, 2021.
\newblock doi:10.1017/jsl.2020.9; arXiv:1805.02095.

\bibitem{PakhomovWalshOmega}
F. Pakhomov and J. Walsh.
\newblock Reducing $\omega$-model reflection to iterated syntactic reflection.
\newblock \emph{Journal of Mathematical Logic}, 23(2), 2023.
\newblock doi:10.1142/S0219061322500015; arXiv:2103.12147.

\bibitem{Rathjen}
M. Rathjen.
\newblock Well-ordering principles in proof theory and reverse mathematics.
\newblock arXiv:2010.12453, 2020.

\end{thebibliography}

\end{document}
